% Generated draft: requires mathematical and source review.
\documentclass[11pt]{article}
\usepackage[margin=25mm]{geometry}
\usepackage{fontspec}
\setmainfont{Latin Modern Roman}
\usepackage{amsmath,amssymb,mathtools,unicode-math}
\setmathfont{Latin Modern Math}
\usepackage{enumitem,needspace,xurl,hyperref}
\setlength{\parindent}{0pt}
\setlength{\parskip}{5pt}
\setlength{\emergencystretch}{3em}
\newcommand{\N}{\mathbb{N}}
\newcommand{\Z}{\mathbb{Z}}
\newcommand{\Q}{\mathbb{Q}}
\newcommand{\R}{\mathbb{R}}
\newcommand{\C}{\mathbb{C}}
\newcommand{\F}{\mathbb{F}}
\newcommand{\A}{\mathbb{A}}
\newcommand{\PP}{\mathbb{P}}
\newcommand{\E}{\mathbb{E}}
\newcommand{\eps}{\varepsilon}
\newcommand{\dd}{\,\mathrm{d}}
\begin{document}
\section*{UnsolvedMath ID: 10400079}
\textbf{Source problem code:} catalogue entry
\subsection{Definitions and mathematical statement}
Problem 4.1 asks to calculate \(\mathcal S_{2,\infty}(M)\) for every oriented \(3\)-manifold \(M\), and to find a convenient general method of calculation. The record omits the preceding definition of this module, including its coefficient ring and parameter.

In the usual Kauffman-bracket skein interpretation of this notation, fix a commutative ring \(R\) with unit and an invertible \(A\in R\). Take the free \(R\)-module on isotopy classes of unoriented framed links in \(M\), including the empty link, and impose the local skein and trivial-circle relations
\[
 [L_{\times}]=A[L_0]+A^{-1}[L_{\infty}],\qquad
 [L\sqcup O]=-(A^2+A^{-2})[L],
\]
where the three links in the first relation agree outside a ball and differ inside it by a crossing and its two prescribed smoothings, and \(O\) is a trivial framed circle in a ball disjoint from \(L\). The resulting module is commonly written \(\mathcal S_{2,\infty}(M;R,A)\). The smoothings and framing conventions require the source's diagrams. This standard interpretation is conditional: the supplied question does not specify \(R\), \(A\), or the precise relations.

After recovering those choices, the requested computation should determine the module for each oriented \(M\), for example by an explicit module presentation or an equivalent structural description, and provide a usable method for obtaining it. No compactness, closedness, boundary condition, input representation, algorithmic complexity target, or permitted output form is stated. In particular the universal quantifier over oriented \(3\)-manifolds must not be silently replaced by a restricted familiar family.
\subsection{Short English statement}
Compute the source's \(\mathcal S_{2,\infty}\) skein module for every oriented three-manifold and develop a convenient way to do so. The usual Kauffman-bracket model explains the notation, but the source's coefficients, parameter and manifold conventions are absent. The supplied summary reports partial computations without naming them.
\subsection{Sources}
\begin{itemize}
\item Local UnsolvedMath registry, permanent ID \(10400079\), exact source code \texttt{catalogue entry}, Problem 4.1. No individual proposer or proposed year is supplied.
\item Ohtsuki, \emph{Problems on invariants of knots and 3-manifolds} (2002), Problem 4.1, PDF page 74: \url{https://msp.org/gtm/2002/04/gtm-2002-04-024s.pdf}. Bold headings distinguish 231 numbered items from references in surrounding exposition.
\item Extraction pdf-font-structure, accessed 2026-07-29. Source evidence NEEDS\_REVIEW presents an open problem and calls for a current primary-literature audit. Supplied OPEN-TRIAGE summary reports partial results for various manifolds but gives no further references.
\item Dataset attribution: Ulam AI / UnsolvedMath Contributors, CC BY 4.0. Source rights NEEDS\_REVIEW: publicly accessible author/publisher copy, redistribution terms requiring review.
\end{itemize}
\end{document}
